\documentclass[../../main.tex]{subfiles}

\begin{document}

\section{预测度的延拓: Carath\'eodory-Hahn 定理}\label{section:ch17-caratheodory-hahn}

令 $\mu:\mathcal{S}\to[0,\infty]$ 为定义在集合 $X$ 的子集的非空族 $\mathcal{S}$ 上的集函数. 我们问以下问题: 为使得由 $\mu$ 诱导的测度 $\overline{\mu}$ 是 $\mu$ 的延拓, 即 $\mathcal{S}$ 中的每个集 $E$ 关于外测度 $\mu^*$ 是可测的, $\mu(E)=\mu^*(E)$, $\mathcal{S}$ 和集函数 $\mu$ 必须具有什么性质? 我们将确定集函数 $\mu$ 必须具有的必要性质, 且证明: 若 $\mathcal{S}$ 有更好的集论结构, 这些性质是充分的.

我们称集函数 $\mu:\mathcal{S}\to[0,\infty]$ 是有限可加的, 若 $\{E_{k}\}_{k=1}^{n}$ 是有限不交集族且 $\bigcup_{k=1}^{n}E_{k}$ 也属于 $\mathcal{S}$, 则
\begin{align*}
\mu\left(\bigcup_{k=1}^{n}E_{k}\right)=\sum_{k=1}^{n}\mu(E_{k}).
\end{align*}

\begin{proposition}\label{proposition:ch17-extension-necessary}
令 $\mathcal{S}$ 为集合 $X$ 的子集族, 而 $\mu:\mathcal{S}\to[0,\infty]$ 是一个集函数. 为使得由 $\mu$ 诱导的 Carath\'eodory 测度是 $\mu$ 的延拓, $\mu$ 必须既是有限可加的又是可数单调的, 且若 $\varnothing$ 属于 $\mathcal{S}$, $\mu(\varnothing)=0$.
\end{proposition}

\begin{proof}
令 $(X,\mathcal{M},\overline{\mu})$ 表示由 $\mu$ 诱导的 Carath\'eodory 测度空间且假定 $\overline{\mu}:\mathcal{M}\to[0,\infty]$ 延拓 $\mu:\mathcal{S}\to[0,\infty]$. 观察到若 $\varnothing$ 属于 $\mathcal{S}$, 则 $\mu(\varnothing)=\overline{\mu}(\varnothing)=0$, 这是由于 $\overline{\mu}$ 是延拓 $\mu$ 的测度.

现在令 $\{E_{k}\}_{k=1}^{n}$ 为 $\mathcal{S}$ 中集合的不交族, 使得 $\bigcup_{k=1}^{n}E_{k}$ 也属于 $\mathcal{S}$. 测度是有限可加的, 因为它是可数可加的且空集的测度为零. 因此, 由于 $\overline{\mu}$ 延拓 $\mu$,
\begin{align*}
\mu\left(\bigcup_{k=1}^{n}E_{k}\right)=\overline{\mu}\left(\bigcup_{k=1}^{n}E_{k}\right)=\sum_{k=1}^{n}\overline{\mu}(E_{k})=\sum_{k=1}^{n}\mu(E_{k}).
\end{align*}
于是 $\mu$ 是有限可加的. 为证明可数单调性, 观察到对所有 $E\in\mathcal{S}$, $\mu(E)=\mu^*(E)$ 当且仅当 $\mu$ 是可数单调的. 于是, 若 $\overline{\mu}$ 延拓 $\mu$, 对所有 $E\in\mathcal{S}$, $\mu^*(E)=\overline{\mu}(E)=\mu(E)$, 因此 $\mu$ 可数单调的.
\end{proof}

该命题暗示着对以下的集函数类单独命名是有用的.

\begin{definition}[预测度]\label{definition:ch17-premeasure}
令 $\mathcal{S}$ 为集合 $X$ 的子集族, 而 $\mu:\mathcal{S}\to[0,\infty]$ 是一个集函数. 则称 $\mu$ 为预测度, 若 $\mu$ 既是有限可加的又是可数单调的, 且若 $\varnothing$ 属于 $\mathcal{S}$, 则 $\mu(\varnothing)=0$.
\end{definition}

对由 $\mu$ 诱导的 Carath\'eodory 测度来说, 成为预测度是它成为 $\mu$ 的延拓的一个必要但不充分的条件. 然而, 若我们在 $\mathcal{S}$ 上施加更好的集论结构, 该必要条件也是充分的.

\begin{definition}[相对补封闭与有限交封闭]\label{definition:ch17-closedness}
$X$ 的子集族 $\mathcal{S}$ 称为关于相对补封闭, 若 $A$ 与 $B$ 属于 $\mathcal{S}$, 相对补 $A\sim B$ 属于 $\mathcal{S}$. 族 $\mathcal{S}$ 称为关于有限交封闭, 若 $A$ 与 $B$ 属于 $\mathcal{S}$, 交 $A\cap B$ 属于 $\mathcal{S}$.
\end{definition}

观察到若子集族 $\mathcal{S}$ 关于相对补封闭, 则它关于有限交封闭, 这是由于若 $A$ 与 $B$ 属于 $\mathcal{S}$, 则
\begin{align*}
A\cap B=A\sim[A\sim B]
\end{align*}
也如此.

也观察到若非空集族 $\mathcal{S}$ 关于相对补封闭, 则它包含 $\varnothing$. 的确, $\varnothing=A\sim A$, 其中 $A$ 属于 $\mathcal{S}$.

\begin{theorem}\label{theorem:ch17-premeasure-extension}
令 $\mu:\mathcal{S}\to[0,\infty]$ 为 $X$ 的子集的非空族 $\mathcal{S}$ 上的预测度. 它关于相对补封闭. 则由 $\mu$ 诱导的 Carath\'eodory 测度 $\overline{\mu}:\mathcal{M}\to[0,\infty]$ 是 $\mu$ 的延拓: 它称为 $\mu$ 的 Carath\'eodory 延拓.
\end{theorem}

\begin{proof}
令 $A$ 属于 $\mathcal{S}$. 为证明 $A$ 关于由 $\mu$ 诱导的外测度可测, 仅须令 $E$ 为 $X$ 的任何有限外测度的子集, 令 $\varepsilon>0$ 并证明
\begin{align}
\mu^*(E)+\varepsilon\geqslant\mu^*(E\cap A)+\mu^*(E\cap A^{c}).\label{align:ch17-proof-ineq1}
\end{align}
根据外测度的定义, 存在覆盖 $E$ 的 $\mathcal{S}$ 中的集族 $\{E_{k}\}_{k=1}^{\infty}$ 满足
\begin{align}
\mu^*(E)+\varepsilon\geqslant\sum_{k=1}^{\infty}\mu(E_{k}).\label{align:ch17-proof-ineq2}
\end{align}
然而, 对每个 $k$, 由于 $\mathcal{S}$ 关于相对补封闭, $E_{k}\sim A$ 属于 $\mathcal{S}$ 且 $E_{k}\cap A=E_{k}\sim[E_{k}\sim A]$ 也如此. 预测度是有限可加的. 因此
\begin{align*}
\mu(E_{k})=\mu(E_{k}\cap A)+\mu(E_{k}\cap A^{c}).
\end{align*}
对这些不等式求和得到
\begin{align}
\sum_{k=1}^{\infty}\mu(E_{k})=\sum_{k=1}^{\infty}\mu(E_{k}\cap A)+\sum_{k=1}^{\infty}\mu(E_{k}\cap A^{c}).\label{align:ch17-proof-ineq3}
\end{align}
观察到 $\{E_{k}\cap A\}_{k=1}^{\infty}$ 和 $\{E_{k}\cap A^{c}\}_{k=1}^{\infty}$ 分别是覆盖 $E\cap A$ 与 $E\cap A^{c}$ 的 $\mathcal{S}$ 中集合的可数族. 因此, 根据外测度的定义,
\begin{align*}
\sum_{k=1}^{\infty}\mu(E_{k}\cap A)\geqslant\mu^*(E\cap A)\text{ 且 }\sum_{k=1}^{\infty}\mu(E_{k}\cap A^{c})\geqslant\mu^*(E\cap A^{c}).
\end{align*}
所需要的不等式\eqref{align:ch17-proof-ineq1}从这两个不等式以及\eqref{align:ch17-proof-ineq2}和\eqref{align:ch17-proof-ineq3}得出.

显然 $\mu(E)=\mu^*(E)$ 对每个集合 $E\in\mathcal{S}$ 成立当且仅当 $\mu$ 是可数单调的. 因此对每个 $E\in\mathcal{S}$, $\mu(E)=\mu^*(E)$, 由于每个集合 $E\in\mathcal{S}$ 是可测的, $\mu(E)=\overline{\mu}(E)$.
\end{proof}

\begin{remark}
观察到在上述定理的证明中预测度的两个性质扮演了十分不同的角色. 我们用 $\mu$ 的有限可加性推出 $\mathcal{S}$ 中的每个集合是 $\mu^*$ 可测的. $\mu$ 的可数单调性等价于等式 $\mu(E)=\mu^*(E)$ 对所有 $E\in\mathcal{S}$ 成立.
\end{remark}

\begin{definition}[半环]\label{definition:ch17-semiring}
集合 $X$ 的子集的非空族 $\mathcal{S}$ 称为半环, 若只要 $A$ 和 $B$ 属于 $\mathcal{S}$, $A\cap B$ 就属于 $\mathcal{S}$, 且存在 $\mathcal{S}$ 中的集合的有限不交族 $\{C_{k}\}_{k=1}^{n}$ 使得
\begin{align*}
A\sim B=\bigcup_{k=1}^{n}C_{k}.
\end{align*}
\end{definition}

\begin{proposition}\label{proposition:ch17-semiring-extension}
令 $\mathcal{S}$ 为集合 $X$ 的子集的半环. 定义 $\mathcal{S}'$ 为 $\mathcal{S}$ 中的集合的有限不交族的并. 则 $\mathcal{S}'$ 关于相对补封闭. 此外, $\mathcal{S}$ 上的任何预测度可唯一地延拓为 $\mathcal{S}'$ 上的预测度.
\end{proposition}

\begin{proof}
显然 $\mathcal{S}'$ 关于有限并与有限交封闭. 令 $\{A_{k}\}_{k=1}^{n}$ 和 $\{B_{j}\}_{j=1}^{m}$ 为 $\mathcal{S}$ 中的集合的两个有限不交族. 观察到
\begin{align}
\left[\bigcup_{k=1}^{n}A_{k}\right]\sim\left[\bigcup_{j=1}^{m}B_{j}\right]=\bigcup_{k=1}^{n}\left[\bigcap_{j=1}^{m}(A_{k}\sim B_{j})\right].\label{align:ch17-semiring-diff}
\end{align}
由于每个 $A_{k}\sim B_{j}$ 属于 $\mathcal{S}'$ 且 $\mathcal{S}'$ 关于有限并与有限交封闭, 我们从\eqref{align:ch17-semiring-diff}推出 $\mathcal{S}'$ 关于相对补封闭.

令 $\mu:\mathcal{S}\to[0,\infty]$ 为 $\mathcal{S}$ 上的预测度. 对满足 $E=\bigcup_{k=1}^{n}A_{k}$ 的 $E\subseteq X$, 其中 $\{A_{k}\}_{k=1}^{n}$ 是 $\mathcal{S}$ 中集合的不交族, 定义 $\mu'(E)=\sum_{k=1}^{n}\mu(A_{k})$. 为证明 $\mu'(E)$ 是恰当定义的, 令 $E$ 也是 $\mathcal{S}$ 中集合的有限族 $\{B_{j}\}_{j=1}^{m}$ 的不交并. 我们必须证明
\begin{align*}
\sum_{j=1}^{m}\mu(B_{j})=\sum_{k=1}^{n}\mu(A_{k}).
\end{align*}
然而, 根据预测度的有限可加性,
\begin{align*}
\mu(B_{j})=\sum_{k=1}^{n}\mu(B_{j}\cap A_{k}),\quad 1\leqslant j\leqslant m
\end{align*}
且
\begin{align*}
\mu(A_{k})=\sum_{j=1}^{m}\mu(B_{j}\cap A_{k}),\quad 1\leqslant k\leqslant n.
\end{align*}
因此
\begin{align*}
\sum_{j=1}^{m}\mu(B_{j})=\sum_{j=1}^{m}\left[\sum_{k=1}^{n}\mu(B_{j}\cap A_{k})\right]=\sum_{k=1}^{n}\left[\sum_{j=1}^{m}\mu(B_{j}\cap A_{k})\right]=\sum_{k=1}^{n}\mu(A_{k}).
\end{align*}
因此 $\mu'$ 在 $\mathcal{S}'$ 上是恰当定义的.

剩下来要证明 $\mu'$ 是 $\mathcal{S}'$ 上的预测度. 由于 $\mu'$ 是恰当定义的, 它继承了 $\mu$ 的有限可加性. 为证明 $\mu'$ 的可数单调性, 令 $E\in\mathcal{S}'$ 被 $\mathcal{S}'$ 中的集族 $\{E_{k}\}_{k=1}^{\infty}$ 覆盖. 不失一般性我们可以假设 $\{E_{k}\}_{k=1}^{\infty}$ 是 $\mathcal{S}'$ 中集合的不交族(见习题 31 的(iii)部分). 令 $E=\bigcup_{j=1}^{m}A_{j}$, 其中并是不交的且每个 $A_{j}$ 属于 $\mathcal{S}$. 对每个 $j$, $A_{j}$ 可被 $\mathcal{S}$ 中集合的可数族 $\bigcup_{k=1}^{\infty}(A_{j}\cap E_{k})$ 覆盖, 根据 $\mu$ 的可数单调性,
\begin{align*}
\mu(A_{j})\leqslant\sum_{k=1}^{\infty}\mu(A_{j}\cap E_{k}).
\end{align*}
因此根据 $\mu$ 的有限单调性,
\begin{align*}
\mu'(E)=\sum_{j=1}^{m}\mu(A_{j})\leqslant\sum_{j=1}^{m}\left[\sum_{k=1}^{\infty}\mu(A_{j}\cap E_{k})\right]=\sum_{k=1}^{\infty}\left[\sum_{j=1}^{m}\mu(A_{j}\cap E_{k})\right]=\sum_{k=1}^{\infty}\mu(E\cap E_{k})\leqslant\sum_{k=1}^{\infty}\mu'(E_{k}).
\end{align*}
因此 $\mu'$ 是可数单调的. 证明完毕.
\end{proof}

对于 $X$ 的子集族 $\mathcal{S}$, 集函数 $\mu:\mathcal{S}\to[0,\infty]$ 称为是 $\sigma$ 有限的, 若 $X=\bigcup_{k=1}^{\infty}S_{k}$, 其中对每个 $k$, $S_{k}\in\mathcal{S}$ 且 $\mu(S_{k})<\infty$.

\begin{theorem}[Carath\'eodory-Hahn]\label{theorem:ch17-caratheodory-hahn}
令 $\mu:\mathcal{S}\to[0,\infty]$ 为 $X$ 的子集的半环 $\mathcal{S}$ 上的预测度. 则由 $\mu$ 诱导的 Carath\'eodory 测度 $\overline{\mu}$ 是 $\mu$ 的延拓. 此外, 若 $\mu$ 是 $\sigma$ 有限的, 则 $\overline{\mu}$ 也如此, 且 $\overline{\mu}$ 是 $\mu^*$ 可测集的 $\sigma$ 代数上的延拓 $\mu$ 的唯一测度.
\end{theorem}

\begin{proof}
我们从定理\cref{theorem:ch17-premeasure-extension}和命题\cref{proposition:ch17-semiring-extension}推出 $\overline{\mu}$ 延拓 $\mu$. 现在假设 $\mu$ 是 $\sigma$ 有限的. 为证明唯一性, 令 $\mu_{1}$ 为 $\mathcal{M}$ 上的另一个延拓 $\mu$ 的测度. 我们把 $X$ 表示为 $X=\bigcup_{k=1}^{\infty}X_{k}$, 其中并是不交的, 且对每个 $k$, $X_{k}$ 属于 $\mathcal{S}$ 且 $\mu(X_{k})<\infty$. 根据测度的可数可加性, 为证明唯一性仅须证明 $\overline{\mu}$ 和 $\mu_{1}$ 在包含于每个 $X_{k}$ 的可测集上一致. 令 $E$ 为可测集且满足 $E\subseteq E_{0}$, 其中 $E_{0}\in\mathcal{S}$ 且 $\mu(E_{0})<\infty$. 我们将证明
\begin{align}
\overline{\mu}(E)=\mu_{1}(E).\label{align:ch17-uniqueness}
\end{align}
根据命题\cref{proposition:ch17-gdelta-approximation}, 存在集合 $A\in\mathcal{S}_{\sigma\delta}$ 使得 $E\subseteq A$ 且 $\overline{\mu}(A\sim E)=0$. 可以假设 $A\subseteq E_{0}$. 然而, 根据 $\mu_{1}$ 的可数单调性, 若 $B$ 是可测的且 $\mu^*(B)=0$, 则 $\mu_{1}(B)=0$. 因此 $\mu_{1}(A\sim E)=0$. 另一方面, 根据 $\mu_{1}$ 与 $\overline{\mu}$ 的可数可加性, 这些测度在 $\mathcal{S}_{\sigma}$ 上一致, 因此根据测度的连续性, 它们在 $E_{0}$ 的属于 $\mathcal{S}_{\sigma\delta}$ 的子集上一致. 因此 $\mu_{1}(A)=\overline{\mu}(A)$. 于是
\begin{align*}
\mu_{1}(A\sim E)=\overline{\mu}(A\sim E),\quad\mu_{1}(A)=\overline{\mu}(A),
\end{align*}
从而\eqref{align:ch17-uniqueness}得证.
\end{proof}

\begin{corollary}\label{corollary:ch17-uniqueness-semiring}
令 $\mathcal{S}$ 为集合 $X$ 的子集族的半环, 而 $\mathcal{B}$ 是包含 $\mathcal{S}$ 的 $X$ 的子集的最小 $\sigma$ 代数. 则 $\mathcal{B}$ 上的两个 $\sigma$ 有限测度相等当且仅当它们在 $\mathcal{S}$ 中的集合上相等.
\end{corollary}

我们在 Carath\'eodory-Hahn 定理中施加在 $\mathcal{S}$ 上的集论限制在一些重要情形中得到满足. 例如, 实数的有界区间族以及平面 $\mathbf{R}^{2}$ 的那些由实数的有界区间的笛卡儿积构成的子集族是半环(见习题 33). 此外, $\mathbf{R}^{n}$ 中的有界区间族是半环. 这允许我们用 Carath\'eodory 的方法构造 $\mathbf{R}^{n}$ 上的 Lebesgue 测度.

我们注意到 Carath\'eodory-Hahn 定理中的唯一性断言可能不成立, 若预测度不假设为 $\sigma$ 有限的(见习题 32).

对读者来说, 熟悉一些与集合 $X$ 的子集族 $\mathcal{S}$ 的性质有关的词汇是有用的. 集族 $\mathcal{S}$ 称为集合的环, 若它关于有限并与相对补封闭, 因此关于有限交封闭. 包含 $X$ 的环称为代数, 而包含 $X$ 的半环称为半代数.

\begin{exercise}\label{exercise:ch17-25}
令 $X$ 为多于一点的集合, 而 $A$ 是 $X$ 的非空真子集. 定义 $\mathcal{S}=\{A,X\}$ 与集函数 $\mu:\mathcal{S}\to[0,\infty]$ 为 $\mu(A)=1$ 且 $\mu(X)=2$. 证明 $\mu:\mathcal{S}\to[0,\infty]$ 是一个预测度. $\mu$ 能延拓为测度吗? $X$ 的关于由 $\mu$ 诱导的外测度 $\mu^*$ 可测的子集是什么?
\end{exercise}

\begin{exercise}\label{exercise:ch17-26}
考虑 $\mathbf{R}$ 的子集族 $\mathcal{S}=\{\varnothing,[0,1],[0,3],[2,3]\}$, 定义 $\mu(\varnothing)=0$, $\mu([0,1])=1$, $\mu([0,3])=1$, $\mu([2,3])=1$. 证明 $\mu:\mathcal{S}\to[0,\infty]$ 是一个预测度. $\mu$ 能延拓为测度吗? $\mathbf{R}$ 的什么子集关于由 $\mu$ 诱导的外测度 $\mu^*$ 是可测的?
\end{exercise}

\begin{exercise}\label{exercise:ch17-27}
令 $\mathcal{S}$ 为集合 $X$ 的子集族, 而 $\mu:\mathcal{S}\to[0,\infty]$ 是集函数. 证明 $\mu$ 是可数单调的当且仅当 $\mu^*$ 是 $\mu$ 的延拓.
\end{exercise}

\begin{exercise}\label{exercise:ch17-28}
证明集函数是预测度, 若它可延拓为测度.
\end{exercise}

\begin{exercise}\label{exercise:ch17-29}
证明 $\sigma$ 代数上的集函数是测度当且仅当它是预测度.
\end{exercise}

\begin{exercise}\label{exercise:ch17-30}
令 $\mathcal{S}$ 为关于有限并与有限交封闭的集族.
\begin{description}
\item[（1）] 证明 $\mathcal{S}_{\sigma}$ 关于可数并与有限交封闭.
\item[（2）] 证明 $\mathcal{S}_{\sigma\delta}$ 中的每个集是 $\mathcal{S}_{\sigma}$ 集的递减序列的交.
\end{description}
\end{exercise}

\begin{exercise}\label{exercise:ch17-31}
令 $\mathcal{S}$ 为集合 $X$ 的子集的半代数, 而 $\mathcal{S}'$ 为 $\mathcal{S}$ 中集合的有限不交族的并所构成的族.
\begin{description}
\item[（i）] 证明 $\mathcal{S}'$ 是代数.
\item[（ii）] 证明 $\mathcal{S}_{\sigma}=\mathcal{S}'_{\sigma}$, 因此 $\mathcal{S}_{\sigma\delta}=\mathcal{S}'_{\sigma\delta}$.
\item[（iii）] 令 $\{E'_{k}\}_{k=1}^{\infty}$ 为 $\mathcal{S}'$ 的集族. 证明我们能将 $\bigcup_{k=1}^{\infty}E'_{k}$ 表示为 $\mathcal{S}$ 中集合的不交并 $\bigcup_{k=1}^{\infty}E_{k}$, 使得
\begin{align*}
\sum_{k=1}^{\infty}\mu'(E'_{k})\geqslant\sum_{k=1}^{\infty}\mu(E_{k}).
\end{align*}
\item[（iv）] 令 $A$ 属于 $\mathcal{S}'_{\sigma\delta}$. 证明 $A$ 是 $\mathcal{S}_{\sigma}$ 中集合的递减序列 $\{A_{k}\}_{k=1}^{\infty}$ 的交.
\end{description}
\end{exercise}

\begin{exercise}\label{exercise:ch17-32}
令 $\mathbf{Q}$ 为有理数集, 而 $\mathcal{S}$ 为形如 $(a,b]\cap\mathbf{Q}$ 的区间的全体有限并, 其中 $a,b\in\mathbf{Q}$ 且 $a\leqslant b$. 若 $a<b$, 定义 $\mu(a,b]=\infty$ 与 $\mu(\varnothing)=0$. 证明 $\mathcal{S}$ 关于相对补是封闭的, 而 $\mu:\mathcal{S}\to[0,\infty]$ 是一个预测度. 接着证明 $\mu$ 到包含 $\mathcal{S}$ 的最小 $\sigma$ 代数的延拓不是唯一的.
\end{exercise}

\begin{exercise}\label{exercise:ch17-33}
谈到实数的有界区间, 我们指的是对实数 $a\leqslant b$ 得到的形如 $[a,b]$, $[a,b)$, $(a,b]$ 或 $(a,b)$ 的集合. 因此考虑空集与单点集为有界区间. 证明以下三个集族 $\mathcal{S}$ 是半环.
\begin{description}
\item[（i）] 令 $\mathcal{S}$ 为实数的所有有界区间构成的族.
\item[（ii）] 令 $\mathcal{S}$ 为 $\mathbf{R}\times\mathbf{R}$ 的由所有实数的有界区间的乘积构成的子集族.
\item[（iii）] 令 $n$ 为自然数, 而 $X$ 是 $\mathbf{R}$ 的 $n$ 重笛卡儿积: $X=\underbrace{\mathbf{R}\times\cdots\times\mathbf{R}}_{n\text{ 次}}$. 令 $\mathcal{S}$ 为 $X$ 的由所有实数的有界区间的 $n$ 重笛卡儿积构成的子集族.
\end{description}
\end{exercise}

\begin{exercise}\label{exercise:ch17-34}
若我们从 $2^{X}$ 上的外测度 $\mu^*$ 开始在 $\mu^*$ 可测集上构造诱导测度 $\overline{\mu}$, 能将 $\overline{\mu}$ 视为集函数且将由 $\overline{\mu}$ 诱导的外测度记为 $\mu^{+}$.
\begin{description}
\item[（i）] 证明对每个集合 $E\subseteq X$, 我们有 $\mu^{+}(E)\geqslant\mu^*(E)$.
\item[（ii）] 对给定集合 $E$, 证明 $\mu^{+}(E)=\mu^*(E)$ 当且仅当存在 $\mu^*$ 可测集 $A\supseteq E$ 满足 $\mu^*(A)=\mu^*(E)$.
\end{description}
\end{exercise}

\begin{exercise}\label{exercise:ch17-35}
令 $\mathcal{S}$ 为 $X$ 的子集的 $\sigma$ 代数, 而 $\mu:\mathcal{S}\to[0,\infty]$ 是测度. 令 $\overline{\mu}:\mathcal{M}\to[0,\infty]$ 为通过 Carath\'eodory 构造法由 $\mu$ 诱导的测度. 证明 $\mathcal{S}$ 是 $\mathcal{M}$ 的子族且它可以是真子族.
\end{exercise}

\begin{exercise}\label{exercise:ch17-36}
令 $\mu$ 为代数 $\mathcal{S}$ 上的有限预测度, 而 $\mu^*$ 是诱导的外测度. 证明 $X$ 的子集 $E$ 是 $\mu^*$ 可测的当且仅当对每个 $\varepsilon>0$, 存在集合 $A\in\mathcal{S}_{\sigma}$, $A\subseteq E$ 使得 $\mu^*(E\sim A)<\varepsilon$.
\end{exercise}

\end{document}